\section{Experiments}
% DRAFT-ONLY: contains artificial placeholders for planning and layout review.

\subsection{Setup}
Table~\ref{tab:setup} reports environments, metrics, training budget, seeds, baselines.
% Auto-generated provisional table. Replace artificial/estimated entries before submission.
\begin{table}[t]
\centering
\small
\begin{tabular}{lcccc}
\toprule
Environment/Task & Primary metric & Training budget & Seeds & Baselines \\
\midrule
Walker Walk & Return (↑) & 500k steps & \textbf{5.00} & TD-MPC2, DreamerV3, PlaNet \\
Cheetah Run & Return (↑) & 500k steps & \underline{5.00} & TD-MPC2, DreamerV3, PlaNet \\
Cartpole Swingup & Return (↑) & 500k steps & 5.00 & TD-MPC2, DreamerV3, PlaNet \\
Reacher Easy & Return (↑) & 500k steps & 5.00 & TD-MPC2, DreamerV3, PlaNet \\
\bottomrule
\end{tabular}
\caption{Setup. Provenance badges mark provisional entries.}
\label{tab:setup}
\end{table}

\subsection{Main Results}
Table~\ref{tab:main_results} reports methods × tasks final-score table.
% Auto-generated provisional table. Replace artificial/estimated entries before submission.
\begin{table}[t]
\centering
\small
\begin{tabular}{lcccc}
\toprule
Method & Walker Walk Return ($\uparrow$) & Cheetah Run Return ($\uparrow$) & Cartpole Swingup Return ($\uparrow$) & Reacher Easy Return ($\uparrow$) \\
\midrule
\textbf{TRACE} & \textbf{840.00} & \textbf{903.00 $\pm$ 16.00}\textcolor{gray}{\scriptsize[ARTIFICIAL]} & \textbf{791.00 $\pm$ 18.00}\textcolor{gray}{\scriptsize[ARTIFICIAL]} & \textbf{954.00 $\pm$ 20.00}\textcolor{gray}{\scriptsize[ARTIFICIAL]} \\
TD-MPC2 & \underline{828.00 $\pm$ 14.00} & \underline{888.00 $\pm$ 16.00} & \underline{773.00 $\pm$ 18.00}\textcolor{gray}{\scriptsize[ARTIFICIAL]} & \underline{933.00 $\pm$ 20.00}\textcolor{gray}{\scriptsize[ARTIFICIAL]} \\
DreamerV3 & 820.00 $\pm$ 14.00 & 880.00 $\pm$ 16.00 & 765.00 $\pm$ 18.00\textcolor{gray}{\scriptsize[ARTIFICIAL]} & 925.00 $\pm$ 20.00\textcolor{gray}{\scriptsize[ARTIFICIAL]} \\
PlaNet & 778.00 $\pm$ 14.00\textcolor{gray}{\scriptsize[ARTIFICIAL]} & 838.00 $\pm$ 16.00\textcolor{gray}{\scriptsize[ARTIFICIAL]} & 723.00 $\pm$ 18.00\textcolor{gray}{\scriptsize[ARTIFICIAL]} & 883.00 $\pm$ 20.00\textcolor{gray}{\scriptsize[ARTIFICIAL]} \\
\bottomrule
\end{tabular}
\caption{Main Results. Provenance badges mark provisional entries.}
\label{tab:main_results}
\end{table}

\subsection{Learning and Sample Efficiency}
Figure~\ref{fig:learning_efficiency} should be exported from the editable workbench as a multi-panel visualization. The draft workbench currently contains 4 panels for one learning-curve panel per task/environment.
% TODO: select the subset of panels to include in the final paper and export finalized SVG/PDF.
\begin{figure}[t]
\centering
\fbox{\parbox{0.9\linewidth}{Placeholder for multi-panel figure: Learning and Sample Efficiency}}
\caption{Learning and Sample Efficiency. Multi-panel provisional visualization; replace with exported figure after final experiments.}
\label{fig:learning_efficiency}
\end{figure}

\subsection{Ablation Study}
Table~\ref{tab:ablation} reports full method and variants.
% Auto-generated provisional table. Replace artificial/estimated entries before submission.
\begin{table}[t]
\centering
\small
\begin{tabular}{lcccc}
\toprule
Variant & Walker Walk Return ($\uparrow$) & Cheetah Run Return ($\uparrow$) & Cartpole Swingup Return ($\uparrow$) & Reacher Easy Return ($\uparrow$) \\
\midrule
\textbf{TRACE} & \textbf{840.00 $\pm$ 10.00} & \textbf{903.00 $\pm$ 10.00}\textcolor{gray}{\scriptsize[ARTIFICIAL]} & \textbf{791.00 $\pm$ 10.00}\textcolor{gray}{\scriptsize[ARTIFICIAL]} & \textbf{954.00 $\pm$ 10.00}\textcolor{gray}{\scriptsize[ARTIFICIAL]} \\
w/o adaptive coding entropy & \underline{828.00 $\pm$ 11.00}\textcolor{gray}{\scriptsize[ARTIFICIAL]} & \underline{887.00 $\pm$ 11.00}\textcolor{gray}{\scriptsize[ARTIFICIAL]} & \underline{771.00 $\pm$ 11.00}\textcolor{gray}{\scriptsize[ARTIFICIAL]} & \underline{930.00 $\pm$ 11.00}\textcolor{gray}{\scriptsize[ARTIFICIAL]} \\
w/o transition-graph regularization & 816.00 $\pm$ 12.00\textcolor{gray}{\scriptsize[ARTIFICIAL]} & 875.00 $\pm$ 12.00\textcolor{gray}{\scriptsize[ARTIFICIAL]} & 759.00 $\pm$ 12.00\textcolor{gray}{\scriptsize[ARTIFICIAL]} & 918.00 $\pm$ 12.00\textcolor{gray}{\scriptsize[ARTIFICIAL]} \\
w/o soft coding tree & 804.00 $\pm$ 13.00\textcolor{gray}{\scriptsize[ARTIFICIAL]} & 863.00 $\pm$ 13.00\textcolor{gray}{\scriptsize[ARTIFICIAL]} & 747.00 $\pm$ 13.00\textcolor{gray}{\scriptsize[ARTIFICIAL]} & 906.00 $\pm$ 13.00\textcolor{gray}{\scriptsize[ARTIFICIAL]} \\
w/o latent graph feedback & 792.00 $\pm$ 14.00\textcolor{gray}{\scriptsize[ARTIFICIAL]} & 851.00 $\pm$ 14.00\textcolor{gray}{\scriptsize[ARTIFICIAL]} & 735.00 $\pm$ 14.00\textcolor{gray}{\scriptsize[ARTIFICIAL]} & 894.00 $\pm$ 14.00\textcolor{gray}{\scriptsize[ARTIFICIAL]} \\
\bottomrule
\end{tabular}
\caption{Ablation Study. Provenance badges mark provisional entries.}
\label{tab:ablation}
\end{table}

\subsection{Generalization}
Figure~\ref{fig:generalization} should be exported from the editable workbench as a multi-panel visualization. The draft workbench currently contains 3 panels for one panel per task showing unseen-condition difficulty curves.
% TODO: select the subset of panels to include in the final paper and export finalized SVG/PDF.
\begin{figure}[t]
\centering
\fbox{\parbox{0.9\linewidth}{Placeholder for multi-panel figure: Generalization}}
\caption{Generalization. Multi-panel provisional visualization; replace with exported figure after final experiments.}
\label{fig:generalization}
\end{figure}

\subsection{Robustness}
Figure~\ref{fig:robustness} should be exported from the editable workbench as a multi-panel visualization. The draft workbench currently contains 3 panels for one panel per perturbation type or task.
% TODO: select the subset of panels to include in the final paper and export finalized SVG/PDF.
\begin{figure}[t]
\centering
\fbox{\parbox{0.9\linewidth}{Placeholder for multi-panel figure: Robustness}}
\caption{Robustness. Multi-panel provisional visualization; replace with exported figure after final experiments.}
\label{fig:robustness}
\end{figure}

\paragraph{Scalability: performance curves}
Figure~\ref{fig:scalability_curve} should be exported from the editable workbench as a multi-panel visualization. The draft workbench currently contains 3 panels for one panel per problem-size axis.
% TODO: select the subset of panels to include in the final paper and export finalized SVG/PDF.
\begin{figure}[t]
\centering
\fbox{\parbox{0.9\linewidth}{Placeholder for multi-panel figure: Scalability: performance curves}}
\caption{Scalability: performance curves. Multi-panel provisional visualization; replace with exported figure after final experiments.}
\label{fig:scalability_curve}
\end{figure}

\paragraph{Scalability and Efficiency: cost table}
Table~\ref{tab:efficiency_table} reports training time, inference time, memory, parameters.
% Auto-generated provisional table. Replace artificial/estimated entries before submission.
\begin{table}[t]
\centering
\small
\begin{tabular}{lcccc}
\toprule
Method & Training time (h) $\downarrow$ & Inference time (ms) $\downarrow$ & Memory (GB) $\downarrow$ & Parameters (M) $\downarrow$ \\
\midrule
\textbf{TRACE} & 7.80 $\pm$ 0.40\textcolor{gray}{\scriptsize[ARTIFICIAL]} & 3.20 $\pm$ 0.20\textcolor{gray}{\scriptsize[ARTIFICIAL]} & 4.80 $\pm$ 0.10\textcolor{gray}{\scriptsize[ARTIFICIAL]} & 18.50\textcolor{gray}{\scriptsize[ARTIFICIAL]} \\
TD-MPC2 & \underline{7.10 $\pm$ 0.40}\textcolor{gray}{\scriptsize[ARTIFICIAL]} & \underline{2.90 $\pm$ 0.20}\textcolor{gray}{\scriptsize[ARTIFICIAL]} & \underline{4.20 $\pm$ 0.10}\textcolor{gray}{\scriptsize[ARTIFICIAL]} & \underline{17.20}\textcolor{gray}{\scriptsize[ARTIFICIAL]} \\
DreamerV3 & 8.40 $\pm$ 0.40\textcolor{gray}{\scriptsize[ARTIFICIAL]} & 3.60 $\pm$ 0.20\textcolor{gray}{\scriptsize[ARTIFICIAL]} & 5.10 $\pm$ 0.10\textcolor{gray}{\scriptsize[ARTIFICIAL]} & 21.10\textcolor{gray}{\scriptsize[ARTIFICIAL]} \\
PlaNet & \textbf{5.60 $\pm$ 0.40}\textcolor{gray}{\scriptsize[ARTIFICIAL]} & \textbf{2.40 $\pm$ 0.20}\textcolor{gray}{\scriptsize[ARTIFICIAL]} & \textbf{3.60 $\pm$ 0.10}\textcolor{gray}{\scriptsize[ARTIFICIAL]} & \textbf{15.40}\textcolor{gray}{\scriptsize[ARTIFICIAL]} \\
\bottomrule
\end{tabular}
\caption{Scalability and Efficiency: cost table. Provenance badges mark provisional entries.}
\label{tab:efficiency_table}
\end{table}

\subsection{Qualitative Analysis}
Figure~\ref{fig:qualitative} should be exported from the editable workbench as a multi-panel visualization. The draft workbench currently contains 3 panels for one qualitative diagnostic panel per task/example.
% TODO: select the subset of panels to include in the final paper and export finalized SVG/PDF.
\begin{figure}[t]
\centering
\fbox{\parbox{0.9\linewidth}{Placeholder for multi-panel figure: Qualitative Analysis}}
\caption{Qualitative Analysis. Multi-panel provisional visualization; replace with exported figure after final experiments.}
\label{fig:qualitative}
\end{figure}

\subsection{Case Study / Failure Cases}
Table~\ref{tab:failure_cases} reports failure modes, frequency, explanation.
% Auto-generated provisional table. Replace artificial/estimated entries before submission.
\begin{table}[t]
\centering
\small
\begin{tabular}{lccc}
\toprule
Failure mode & Frequency (\%) $\downarrow$ & Explanation & Mitigation \\
\midrule
latent module collapse\textcolor{gray}{\scriptsize[ARTIFICIAL]} & \underline{12.00}\textcolor{gray}{\scriptsize[ARTIFICIAL]} & regularization too strong\textcolor{gray}{\scriptsize[ARTIFICIAL]} & reduce lambda\_TRACE or add warmup\textcolor{gray}{\scriptsize[ARTIFICIAL]} \\
poor long-horizon rollout\textcolor{gray}{\scriptsize[ARTIFICIAL]} & 18.00\textcolor{gray}{\scriptsize[ARTIFICIAL]} & world-model compounding error\textcolor{gray}{\scriptsize[ARTIFICIAL]} & shorter imagination horizon or uncertainty penalty\textcolor{gray}{\scriptsize[ARTIFICIAL]} \\
sparse reward sensitivity\textcolor{gray}{\scriptsize[ARTIFICIAL]} & \textbf{9.00}\textcolor{gray}{\scriptsize[ARTIFICIAL]} & weak early transition signal\textcolor{gray}{\scriptsize[ARTIFICIAL]} & pretraining or reward-balanced replay\textcolor{gray}{\scriptsize[ARTIFICIAL]} \\
\bottomrule
\end{tabular}
\caption{Case Study / Failure Cases. Provenance badges mark provisional entries.}
\label{tab:failure_cases}
\end{table}
